% ====================================================================
% R7-F1 그림 초안 — 블렌드 f_Si(wt% 스윕) 가족 dQ/dV 곡선  [V22-R7-F1]
%   라벨 제안 = fig:blend-family / 삽입 후보 = §3.3 eq:blend-dqdv 부근(1순위) 또는 §3.5 코드 예시
%   ★초안 전용: _sections/ 미수정. 이 블록을 그대로 §3.3 keybox 뒤 또는 §3.5 에 \input/붙여넣기.
%   좌표 = Anode_Fit_v1.0.22.py 의 BlendedAnodeDQDV.from_wt(m_Si, si_case='elemental') 실계산
%     (equilibrium(V,298.15 K) = eq:blend-dqdv 의 |I|→0 평형 종; C_bg=0; 총용량 Q 정규화, 면적=1).
%   검증 로그 = results/comp_R7/F1_NOTE.md (피크 위치·높이 수기 대조, wt%→f_Si 환산, bit-exact).
%   시연값 등급: 전이 중심·폭 = tier-C 시연(SI_ELEMENTAL_LIT; 원소 Si 0.2–0.5 V 범위,
%     obrovac_chevrier2014 tier-B), q_Si=1000 mAh/g = tier-A(limthongkul2003). elemental 케이스는
%     '확인 필요' 공백 0건(SiOₓ 절대전위 공백 등은 미사용 — 정직 규율).
% ====================================================================
\begin{figure}[t]
\centering
\begin{tikzpicture}
% ============================================================
% (a) 전체 창 — 총 평형 dQ/dV 를 총용량 Q 로 정규화한 가족 (4곡선)
%     x = V_n [V] (0.03–0.62) , y = (dQ/dV)/Q [1/V] (면적=1)
% ============================================================
\begin{scope}[x=9.5cm,y=0.58cm]
% ---- 축 ----
\draw[-{Latex[length=1.6mm]}] (0.026,0) -- (0.026,7.95)
  node[above,font=\scriptsize] {$(\dd Q/\dd V)/Q$ [1/V]};
\draw[-{Latex[length=1.6mm]}] (0.026,0) -- (0.642,0)
  node[below,font=\scriptsize] {$V_n$ [V]};
\foreach \xx in {0.1,0.2,0.3,0.4,0.5}
  {\draw (\xx,0.09) -- (\xx,-0.09) node[below,font=\scriptsize] {\xx};}
\foreach \yy in {2,4,6}
  {\draw (0.0235,\yy) -- (0.0285,\yy) node[left,font=\scriptsize] {\yy};}
% ---- 가족 곡선 (흑백; 선종+명도로 구분) ----
% 0 wt%  f_Si=0  (흑연 단독, 기준선) — very thick 실선
\draw[very thick] plot[smooth] coordinates {(0.0300,2.2213) (0.0360,2.6753) (0.0420,3.1886) (0.0480,3.7540) (0.0540,4.3580) (0.0600,4.9795) (0.0660,5.5915) (0.0720,6.1625) (0.0780,6.6611) (0.0840,7.0604) (0.0900,7.3411) (0.0960,7.4944) (0.1020,7.5208) (0.1080,7.4279) (0.1140,7.2282) (0.1200,6.9362) (0.1260,6.5677) (0.1320,6.1396) (0.1380,5.6701) (0.1440,5.1785) (0.1500,4.6844) (0.1560,4.2061) (0.1620,3.7593) (0.1680,3.3554) (0.1740,3.0012) (0.1800,2.6980) (0.1860,2.4423) (0.1920,2.2268) (0.1980,2.0416) (0.2040,1.8765) (0.2100,1.7222) (0.2160,1.5717) (0.2220,1.4213) (0.2280,1.2704) (0.2340,1.1206) (0.2400,0.9750) (0.2460,0.8371) (0.2520,0.7097) (0.2580,0.5950) (0.2640,0.4940) (0.2700,0.4066) (0.2760,0.3322) (0.2820,0.2699) (0.2880,0.2181) (0.2940,0.1756) (0.3000,0.1409) (0.3160,0.0774) (0.3320,0.0421) (0.3480,0.0227) (0.3640,0.0122) (0.3800,0.0066) (0.3960,0.0035) (0.4120,0.0019) (0.4280,0.0010) (0.4440,0.0005) (0.4600,0.0003) (0.4760,0.0002) (0.4920,0.0001) (0.6200,0.0000)};
% 10 wt%  f_Si=0.230 — thick densely dashed
\draw[thick,densely dashed] plot[smooth] coordinates {(0.0300,1.7358) (0.0360,2.0880) (0.0420,2.4861) (0.0480,2.9247) (0.0540,3.3932) (0.0600,3.8757) (0.0660,4.3510) (0.0720,4.7953) (0.0780,5.1844) (0.0840,5.4973) (0.0900,5.7196) (0.0960,5.8442) (0.1020,5.8717) (0.1080,5.8081) (0.1140,5.6629) (0.1200,5.4473) (0.1260,5.1737) (0.1320,4.8550) (0.1380,4.5053) (0.1440,4.1396) (0.1500,3.7728) (0.1560,3.4192) (0.1620,3.0909) (0.1680,2.7967) (0.1740,2.5418) (0.1800,2.3272) (0.1860,2.1502) (0.1920,2.0050) (0.1980,1.8842) (0.2040,1.7796) (0.2100,1.6840) (0.2160,1.5919) (0.2220,1.5004) (0.2280,1.4085) (0.2340,1.3176) (0.2400,1.2297) (0.2460,1.1473) (0.2520,1.0723) (0.2580,1.0060) (0.2640,0.9490) (0.2700,0.9012) (0.2760,0.8618) (0.2820,0.8297) (0.2880,0.8040) (0.2940,0.7833) (0.3000,0.7666) (0.3160,0.7349) (0.3320,0.7134) (0.3480,0.6968) (0.3640,0.6845) (0.3800,0.6767) (0.3960,0.6724) (0.4120,0.6677) (0.4280,0.6570) (0.4440,0.6341) (0.4600,0.5954) (0.4760,0.5407) (0.4920,0.4741) (0.5080,0.4019) (0.5240,0.3304) (0.5400,0.2646) (0.5560,0.2074) (0.5720,0.1599) (0.5880,0.1216) (0.6040,0.0916) (0.6200,0.0685)};
% 20 wt%  f_Si=0.402 — thick densely dashdotted, black!75
\draw[thick,densely dashdotted,black!75] plot[smooth] coordinates {(0.0300,1.3729) (0.0360,1.6490) (0.0420,1.9610) (0.0480,2.3047) (0.0540,2.6720) (0.0600,3.0504) (0.0660,3.4237) (0.0720,3.7733) (0.0780,4.0804) (0.0840,4.3288) (0.0900,4.5073) (0.0960,4.6105) (0.1020,4.6389) (0.1080,4.5971) (0.1140,4.4927) (0.1200,4.3343) (0.1260,4.1315) (0.1320,3.8946) (0.1380,3.6345) (0.1440,3.3628) (0.1500,3.0913) (0.1560,2.8310) (0.1620,2.5912) (0.1680,2.3791) (0.1740,2.1984) (0.1800,2.0500) (0.1860,1.9318) (0.1920,1.8393) (0.1980,1.7665) (0.2040,1.7072) (0.2100,1.6555) (0.2160,1.6071) (0.2220,1.5594) (0.2280,1.5118) (0.2340,1.4649) (0.2400,1.4202) (0.2460,1.3792) (0.2520,1.3433) (0.2580,1.3132) (0.2640,1.2893) (0.2700,1.2710) (0.2760,1.2576) (0.2820,1.2483) (0.2880,1.2419) (0.2940,1.2376) (0.3000,1.2343) (0.3160,1.2265) (0.3320,1.2153) (0.3480,1.2008) (0.3640,1.1871) (0.3800,1.1777) (0.3960,1.1724) (0.4120,1.1654) (0.4280,1.1473) (0.4440,1.1078) (0.4600,1.0402) (0.4760,0.9448) (0.4920,0.8285) (0.5080,0.7023) (0.5240,0.5774) (0.5400,0.4624) (0.5560,0.3625) (0.5720,0.2794) (0.5880,0.2125) (0.6040,0.1601) (0.6200,0.1197)};
% 30 wt%  f_Si=0.535 — thick densely dotted, black!58
\draw[thick,densely dotted,black!58] plot[smooth] coordinates {(0.0300,1.0913) (0.0360,1.3084) (0.0420,1.5536) (0.0480,1.8237) (0.0540,2.1124) (0.0600,2.4101) (0.0660,2.7043) (0.0720,2.9803) (0.0780,3.2238) (0.0840,3.4222) (0.0900,3.5667) (0.0960,3.6533) (0.1020,3.6824) (0.1080,3.6576) (0.1140,3.5847) (0.1200,3.4707) (0.1260,3.3230) (0.1320,3.1495) (0.1380,2.9589) (0.1440,2.7602) (0.1500,2.5626) (0.1560,2.3746) (0.1620,2.2036) (0.1680,2.0550) (0.1740,1.9319) (0.1800,1.8349) (0.1860,1.7623) (0.1920,1.7107) (0.1980,1.6752) (0.2040,1.6510) (0.2100,1.6333) (0.2160,1.6188) (0.2220,1.6053) (0.2280,1.5920) (0.2340,1.5792) (0.2400,1.5679) (0.2460,1.5591) (0.2520,1.5535) (0.2580,1.5516) (0.2640,1.5532) (0.2700,1.5579) (0.2760,1.5647) (0.2820,1.5730) (0.2880,1.5817) (0.2940,1.5900) (0.3000,1.5973) (0.3160,1.6079) (0.3320,1.6046) (0.3480,1.5918) (0.3640,1.5770) (0.3800,1.5664) (0.3960,1.5603) (0.4120,1.5516) (0.4280,1.5278) (0.4440,1.4753) (0.4600,1.3854) (0.4760,1.2583) (0.4920,1.1035) (0.5080,0.9354) (0.5240,0.7690) (0.5400,0.6159) (0.5560,0.4828) (0.5720,0.3721) (0.5880,0.2831) (0.6040,0.2132) (0.6200,0.1594)};
% ---- 추세 주석: 흑연 축소(아래) · Si 성장(위) ----
\draw[-{Latex[length=1.4mm]},thick] (0.132,6.35) -- (0.132,3.55);
\node[font=\tiny,anchor=west,align=left] at (0.146,5.75) {흑연 staging 겹침\\envelope 축소 ($m_\mathrm{Si}\!\uparrow$)};
\draw[-{Latex[length=1.4mm]},thick] (0.320,0.25) -- (0.320,1.48);
\node[font=\tiny,anchor=west,align=left] at (0.345,1.92) {Si 기여 성장\\($m_\mathrm{Si}\!\uparrow$)};
% ---- 범례 (우상단 빈 영역) ----
\draw[very thick] (0.400,7.35) -- (0.437,7.35);
\node[font=\tiny,anchor=west] at (0.443,7.35) {$0$ wt\% ($f_\mathrm{Si}{=}0$, 흑연 단독)};
\draw[thick,densely dashed] (0.400,6.65) -- (0.437,6.65);
\node[font=\tiny,anchor=west] at (0.443,6.65) {$10$ wt\% ($f_\mathrm{Si}{=}0.23$)};
\draw[thick,densely dashdotted,black!75] (0.400,5.95) -- (0.437,5.95);
\node[font=\tiny,anchor=west] at (0.443,5.95) {$20$ wt\% ($f_\mathrm{Si}{=}0.40$)};
\draw[thick,densely dotted,black!58] (0.400,5.25) -- (0.437,5.25);
\node[font=\tiny,anchor=west] at (0.443,5.25) {$30$ wt\% ($f_\mathrm{Si}{=}0.54$)};
\node[font=\scriptsize] at (0.33,-1.05) {(a) 전체 창 --- 총 dQ/dV$/Q$};
\end{scope}
% ============================================================
% (b) Si 기여 확대 — host_contributions 의 Si 몫만(eq:blend-dqdv Si 이중합)/Q
%     x = V_n [V] (0.15–0.62) , y = (dQ/dV)_Si/Q [1/V]
% ============================================================
\begin{scope}[xshift=6.4cm,x=9.5cm,y=2.6cm]
\draw[-{Latex[length=1.6mm]}] (0.146,0) -- (0.146,1.74)
  node[above,font=\scriptsize] {$(\dd Q/\dd V)_\mathrm{Si}/Q$ [1/V]};
\draw[-{Latex[length=1.6mm]}] (0.146,0) -- (0.644,0)
  node[below,font=\scriptsize] {$V_n$ [V]};
\foreach \xx in {0.2,0.3,0.4,0.5}
  {\draw (\xx,0.02) -- (\xx,-0.02) node[below,font=\scriptsize] {\xx};}
\foreach \yy in {0.5,1.0,1.5}
  {\draw (0.1445,\yy) -- (0.1475,\yy) node[left,font=\scriptsize] {\yy};}
% 10 wt%
\draw[thick,densely dashed] plot[smooth] coordinates {(0.1500,0.1658) (0.1590,0.1882) (0.1680,0.2130) (0.1770,0.2402) (0.1860,0.2696) (0.1950,0.3012) (0.2040,0.3347) (0.2130,0.3698) (0.2220,0.4059) (0.2310,0.4426) (0.2400,0.4790) (0.2490,0.5143) (0.2580,0.5478) (0.2670,0.5786) (0.2760,0.6059) (0.2850,0.6292) (0.2940,0.6481) (0.3030,0.6624) (0.3120,0.6723) (0.3210,0.6782) (0.3300,0.6808) (0.3390,0.6809) (0.3480,0.6793) (0.3570,0.6770) (0.3660,0.6746) (0.3750,0.6725) (0.3840,0.6711) (0.3930,0.6700) (0.4020,0.6688) (0.4110,0.6666) (0.4200,0.6624) (0.4290,0.6552) (0.4380,0.6439) (0.4470,0.6278) (0.4560,0.6063) (0.4650,0.5797) (0.4740,0.5482) (0.4830,0.5126) (0.4920,0.4741) (0.5010,0.4337) (0.5100,0.3928) (0.5190,0.3523) (0.5280,0.3132) (0.5370,0.2763) (0.5460,0.2421) (0.5550,0.2107) (0.5640,0.1824) (0.5730,0.1572) (0.5820,0.1349) (0.5910,0.1154) (0.6000,0.0984) (0.6090,0.0837) (0.6180,0.0711)};
% 20 wt%
\draw[thick,densely dashdotted,black!75] plot[smooth] coordinates {(0.1500,0.2897) (0.1590,0.3290) (0.1680,0.3723) (0.1770,0.4197) (0.1860,0.4711) (0.1950,0.5263) (0.2040,0.5849) (0.2130,0.6462) (0.2220,0.7094) (0.2310,0.7734) (0.2400,0.8370) (0.2490,0.8988) (0.2580,0.9574) (0.2670,1.0111) (0.2760,1.0589) (0.2850,1.0996) (0.2940,1.1325) (0.3030,1.1575) (0.3120,1.1748) (0.3210,1.1852) (0.3300,1.1897) (0.3390,1.1899) (0.3480,1.1872) (0.3570,1.1831) (0.3660,1.1789) (0.3750,1.1753) (0.3840,1.1728) (0.3930,1.1709) (0.4020,1.1687) (0.4110,1.1649) (0.4200,1.1576) (0.4290,1.1450) (0.4380,1.1253) (0.4470,1.0971) (0.4560,1.0596) (0.4650,1.0130) (0.4740,0.9580) (0.4830,0.8958) (0.4920,0.8285) (0.5010,0.7579) (0.5100,0.6864) (0.5190,0.6157) (0.5280,0.5474) (0.5370,0.4829) (0.5460,0.4230) (0.5550,0.3683) (0.5640,0.3188) (0.5730,0.2748) (0.5820,0.2358) (0.5910,0.2017) (0.6000,0.1720) (0.6090,0.1463) (0.6180,0.1242)};
% 30 wt%
\draw[thick,densely dotted,black!58] plot[smooth] coordinates {(0.1500,0.3859) (0.1590,0.4381) (0.1680,0.4958) (0.1770,0.5590) (0.1860,0.6275) (0.1950,0.7010) (0.2040,0.7790) (0.2130,0.8607) (0.2220,0.9448) (0.2310,1.0301) (0.2400,1.1149) (0.2490,1.1972) (0.2580,1.2751) (0.2670,1.3468) (0.2760,1.4104) (0.2850,1.4646) (0.2940,1.5084) (0.3030,1.5417) (0.3120,1.5648) (0.3210,1.5786) (0.3300,1.5846) (0.3390,1.5848) (0.3480,1.5812) (0.3570,1.5758) (0.3660,1.5701) (0.3750,1.5654) (0.3840,1.5620) (0.3930,1.5595) (0.4020,1.5566) (0.4110,1.5516) (0.4200,1.5419) (0.4290,1.5250) (0.4380,1.4988) (0.4470,1.4612) (0.4560,1.4113) (0.4650,1.3493) (0.4740,1.2759) (0.4830,1.1932) (0.4920,1.1034) (0.5010,1.0095) (0.5100,0.9142) (0.5190,0.8200) (0.5280,0.7291) (0.5370,0.6432) (0.5460,0.5634) (0.5550,0.4905) (0.5640,0.4247) (0.5730,0.3659) (0.5820,0.3141) (0.5910,0.2686) (0.6000,0.2291) (0.6090,0.1948) (0.6180,0.1654)};
% 성장 방향 주석 + Si 창 안내
\draw[-{Latex[length=1.4mm]},thick] (0.240,0.55) -- (0.240,1.42);
\node[font=\tiny,anchor=west] at (0.152,1.62) {Si host 기여항 (eq~\eqref{eq:blend-dqdv} Si 몫); $0$ wt\%: 기여 없음};
\node[font=\tiny,anchor=west] at (0.255,1.30) {$m_\mathrm{Si}\!\uparrow$};
\node[font=\scriptsize] at (0.39,-0.235) {(b) Si 기여 확대};
\end{scope}
\end{tikzpicture}
\caption{블렌드 $f_\mathrm{Si}$(wt\% 스윕) 가족 평형 dQ/dV --- 좌표는 코드
\texttt{BlendedAnodeDQDV.from\_wt}$(m_\mathrm{Si},\,\texttt{si\_case='elemental'})$ 의 실계산이다:
평형식~\eqref{eq:blend-dqdv} 의 $|I|\!\to\!0$ 종
$\sum_{\mathrm{host}}\sum_j Q_j^\mathrm{host}\xi_{\eq,j}^\mathrm{host}(1-\xi_{\eq,j}^\mathrm{host})/w_j^\mathrm{host}$
을 $T{=}298.15$~K$\cdot$$C_\bg{=}0$ 에서 평가하고 총용량 $Q{=}Q_\mathrm{gr}{+}Q_\mathrm{Si}$ 로
정규화했다(곡선 아래 면적 $=1$ 규약, 곡선당 $\sim$$53$--$66$ 표본점). 질량$\leftrightarrow$용량 환산은
§\ref{ssec:si-blend-derive} 각주의 C-052
$f_\mathrm{Si}=m_\mathrm{Si}q_\mathrm{Si}/[m_\mathrm{Si}q_\mathrm{Si}+(1{-}m_\mathrm{Si})q_\mathrm{gr}]$
($q_\mathrm{Si}{=}1000$ mAh/g 원소 Si 1차 가역\cite{limthongkul2003}, tier-A$\cdot$$q_\mathrm{gr}{=}372$)로,
$m_\mathrm{Si}{=}0/10/20/30$ wt\% $\Rightarrow f_\mathrm{Si}{=}0/0.230/0.402/0.535$
($Q{=}0.97/1.26/1.62/2.09\,Q_\cell$). \textbf{(a)} 전체 창: 흑연 staging 4-peak 겹침 envelope 는
$m_\mathrm{Si}$ 증가로 정규화 높이가 $7.5\!\to\!3.7$ [1/V] 로 \emph{축소}하고($Q_\mathrm{gr}$ 절대값은
불변이나 총용량 $Q$ 가 커져 상대 몫 $1{-}f_\mathrm{Si}$ 로 줄어듦), 원소 Si 두 경사 전이(중심 $0.30/0.45$
V, 폭 $0.060/0.050$ V)의 기여가 고전위 쪽에서 \emph{성장}해 $V_n\!\approx\!0.26$ V 부근에서 순서가
역전한다. \textbf{(b)} 그 Si host 몫만 떼어(코드 \texttt{host\_contributions}) 확대: $0\!\to\!1.585$
[1/V] 로 자란다($f_\mathrm{Si}$ 에 정비례). 곡선 개형(전이 중심$\cdot$폭)은 원소 Si 평균 전위대
$0.2$--$0.5$ V\cite{obrovac_chevrier2014}(tier-B) 안의 \textbf{tier-C 시연값}이며(신뢰값 아님, 피팅
override 전제), \texttt{elemental} 케이스는 SiO$_x$ 절대전위 같은 ``확인 필요'' 공백을 쓰지 않는다.
$f_\mathrm{Si}\!\to\!0$ 에서 (a)의 실선은 흑연 단독 평형곡선과 부동소수점까지 일치한다(bit-exact,
식~\eqref{eq:blend-limit} 의 코드판; \S\ref{sec:sum} 그림~\ref{fig:sumcurve} 와 동일 흑연 골격).}
\label{fig:blend-family}
\end{figure}
% 자산(v1.0.22 R7 신규): [V22-R7-F1 블렌드 f_Si wt% 스윕 가족 dQ/dV 2패널 fig:blend-family —
%   eq:blend-dqdv 실계산·elemental·총용량 정규화(면적1)·wt%→f_Si C-052·f_Si→0 bit-exact 검증분]
